% !TEX root = research_plan.tex
%
\usepackage{hyperref}
\usepackage{xcolor}
%
\usepackage[justification=centering]{caption}
%
\usepackage{tikz}
\usetikzlibrary{arrows.meta,positioning,petri,cd}
%
\usepackage[framemethod=TikZ]{mdframed}
\usepackage{amsthm}%\usepackage[framed]{ntheorem}
\usepackage[capitalise]{cleveref}
\mdfdefinestyle{theorembox}{
    linecolor=black,
    linewidth=4pt,
    % --- This controls the "height" of the box around the text ---
    innertopmargin=0pt,     % Space at the top
    innerbottommargin=0pt,  % Space at the bottom
    innerleftmargin=0pt,   % Space on the left
    innerrightmargin=0pt,  % Space on the right
    % ------------------------------------------------------------
    roundcorner=0pt,        % Set to 5pt for rounded
    %backgroundcolor=gray!5, % Light background
    skipabove=0pt,%\baselineskip,
    skipbelow=0pt,%\baselineskip
}
\global\mdfdefinestyle{thm}{%
linecolor=black,
linewidth=1pt,
innertopmargin=0pt
}
\usepackage[framemethod=TikZ]{mdframed}
\usepackage{amsthm}
\usepackage{thmtools}
\usepackage[framemethod=TikZ]{mdframed}
\usepackage{amsthm}
\usepackage{thmtools}
\usepackage{xcolor}

% 1. Define the Frame appearance
\mdfdefinestyle{myThmFrame}{
    innertopmargin=4pt,
    innerbottommargin=4pt,
    innerleftmargin=10pt,
    innerrightmargin=10pt,
    linecolor=black,
    linewidth=1pt,
    %backgroundcolor=blue!5,
    roundcorner=0pt
}

% 2. Define the Theorem Style (The "All-in-one" definition)
\declaretheoremstyle[
    headfont=\bfseries,       % Bold title
    notefont=\itshape,        % Italicized optional note
    bodyfont=\normalfont,     % Standard text for content
    headpunct={. },           % Punctuation after "Theorem 1.1"
    postheadspace=0pt,        % MAGIC: 0pt forces text to stay on the same line
    mdframed={style=myThmFrame} % Attach the frame style here
]{framedrunin}

\declaretheorem[style=framedrunin, name=Theorem]{theorem}%[chapter]
\declaretheorem[style=framedrunin, name=Lemma,      sibling=theorem]{lemma}
\declaretheorem[style=framedrunin, name=Conjecture, sibling=theorem]{conjecture}
\declaretheorem[style=framedrunin, name=Question,   sibling=theorem]{question}
\declaretheorem[style=framedrunin, name=Note,       sibling=theorem]{note}
\declaretheorem[style=framedrunin, name=Example,    sibling=theorem]{example}
\declaretheorem[style=framedrunin, name=Definition, sibling=theorem]{definition}
\declaretheorem[style=framedrunin, name=Problem,    sibling=theorem]{problem}
\declaretheorem[style=framedrunin, name=Problem]{goal}

\newcommand\prob[4]{
  \begin{problem}[#1]\label{#2}
  $ $\\
  {\bfseries Input}: #3\\
  {\bfseries Question}: #4
  \end{problem}
}

%
%%% General text macros %%%
%
\newcounter{commentnumber}
\newcommand{\arka}[1]{\textbf{\textcolor{red}{Arka \stepcounter{commentnumber}\thecommentnumber: #1}}}
\newcommand{\mine}[1]{\textcolor{magenta}{\cite{#1}}}
\newcommand{\mythm}[2]{\textcolor{magenta}{\cite[#1]{#2}}}
\newcommand{\emphdef}[1]{\textcolor{purple}{\emph{#1}}}
\newcommand{\compsci}[1]{\textcolor{blue}{#1}}
\newcommand{\hlt}[1]{\textbf{\textcolor{violet}{#1}}}
%
%%% Domains %%%
%
\newcommand{\Q}{\mathbb{Q}}
\newcommand{\X}{\mathcal{X}}
\newcommand{\M}{\mathcal{M}}
\newcommand{\C}{\mathbb{C}}
\newcommand{\dloQ}{\mathcal{Q}}
\newcommand{\A}{\mathcal{A}}
\newcommand{\B}{\mathcal{B}}
\newcommand{\R}{\mathbb{R}}
\newcommand{\N}{\mathbb{N}}
\newcommand{\Z}{\mathbb{Z}}
\newcommand{\bitV}[2]{{\mathcal{V}_{#1}}{(#2)}}
\newcommand{\ff}{\mathbb{F}}
\newcommand{\fftwo}{\mathbb{F}_2}
\newcommand{\ordZ}{\mathcal{Z}}
\newcommand{\cU}{\mathcal{U}}
%
%%% General math macros
%
\newcommand{\setof}[2]{\left\{#1\ \mid\ #2 \right\}}
\newcommand{\pn}{\mathbf{P}}
%
%%% Atoms %%%
%
\newcommand{\aut}[1]{\mathsf{Aut}{\left(#1\right)}}
\newcommand{\age}[2][]{\textsc{Age}_{#1}{\left(#2\right)}}
\newcommand{\tup}[1]{\overline{#1}}
\newcommand{\otuequiv}[1]{\A^{(#1)}}
%\newcommand{\eq}[1]{\A^{(#1)}}
%
%%% Polynomials %%%
%
\newcommand{\poly}[2]{#1\left[#2\right]}
\newcommand{\idl}[1]{\mathcal{#1}}
\newcommand{\idlgen}[1]{\langle #1 \rangle}
%
%%% vector spaces %%%
%
\newcommand{\flin}[2][\Q]{#2 \to_{\mathsf{fin}} #1}
\newcommand{\bigd}[2][\Q]{#2 \to #1}
\newcommand{\defd}[2][\Q]{#2 \to_{\textsc{def}} #1}
\newcommand{\vgen}[1]{\langle #1 \rangle}
\newcommand{\ofspan}[2]{$\textsc{Span}_{#1}(#2)$}
\newcommand{\idvec}[1]{\mathbf{1}_{#1}}
%
%%% register automata %%%
%
\newcommand{\parikh}{\textsc{Parikh}}
%
%%% complexity %%%
%
\newcommand{\ptime}{\textsc{Ptime}}
\newcommand{\exptime}{\textsc{Exptime}}
\newcommand{\npc}{\textsc{NP-complete}}
